Recommendation has been the largest single deployment of machine learning on the Web for nearly two decades. Until 2022 the field was dominated by one architectural pattern. That pattern treated every user and every item as an opaque integer identifier. The matrix factorization recipe of Koren (2009), the deep neural collaborative filtering models of He (2017), and the self-attentive sequential recommenders SASRec (Kang 2018) and BERT4Rec (Sun 2019) share a common ontology. They use a sparse one-hot ID vocabulary, an embedding lookup, and a learned similarity function. This ID-centric pipeline is efficient at large scale. Production systems at Meta, Alibaba, and TikTok serve catalogs of \(10^9\) items at sub-50 ms latency. Yet the pipeline is inherently closed-world. It cannot reason about the meaning of a movie title. It cannot transfer knowledge across domains. It cannot interact through natural language. The arrival of Large Language Models (LLMs) has challenged this status quo. GPT-3 (Brown 2020), T5 (Raffel 2020), LLaMA (Touvron 2023), and the GPT-4, Claude, and Gemini families package broad open-world knowledge, multi-step reasoning, instruction following, and a unified text interface in a single foundation. The present survey examines how this transformation is unfolding. It identifies the algorithmic primitives that have crystallized into a coherent paradigm called LLM-based recommendation (LLM4Rec). It also charts where the technology, the evaluation methodology, an\ldots{}
